%LTeX: language=fr-FR
\documentclass[a4paper,11pt]{article}

\usepackage[french]{babel}
\usepackage{amsmath,amssymb}
\usepackage{enumitem}
\usepackage{tikz}
\usepackage{tkz-tab}
\usepackage{tasks}
\usetikzlibrary{arrows.meta}
\usepackage[margin=2cm]{geometry}
\usepackage{needspace}
\usepackage{currfile}
\usepackage{fancyhdr}
\usepackage{lastpage}
\usepackage{xcolor}
\settasks{label=\arabic*.}

\pagestyle{empty}
\input{\currfiledir ../def.tex}
\def \Document {Corrigé des exercices}


\setlength{\parindent}{0pt}

\newenvironment{exercice}[1]{%
  \Needspace{6\baselineskip}\par\bigskip\noindent\textbf{Exercice #1}\par\smallskip\noindent\ignorespaces
}{%
  \par
}

% Repère quadrillé : \repere{xmin}{xmax}{ymin}{ymax}
\newcommand{\repere}[4]{%
  \draw[very thin, gray!40] (#1,#3) grid (#2,#4);
  \draw[->] (#1,0) -- ({#2+0.4},0) node[below right] {\small $x$};
  \draw[->] (0,#3) -- (0,{#4+0.4}) node[above left] {\small $y$};
  \foreach \x in {#1,...,#2} {\ifnum\x=0\else\draw (\x,0.08) -- (\x,-0.08) node[below] {\scriptsize $\x$};\fi}
  \foreach \y in {#3,...,#4} {\ifnum\y=0\else\draw (0.08,\y) -- (-0.08,\y) node[left] {\scriptsize $\y$};\fi}
  \node[below left] at (0,0) {\scriptsize $0$};
}

% Tableau de signes : \tabsignes{f}{valeur qui annule}{signes}
\newcommand{\tabsignes}[3]{%
  \begin{tikzpicture}[baseline=(current bounding box.center)]
    \tkzTabInit[lgt=1.2,espcl=1.6]{$x$/0.8,$#1(x)$/0.8}{$-\infty$,$#2$,$+\infty$}
    \tkzTabLine{,#3,}
  \end{tikzpicture}%
}

\begin{document}
\pagestyle{fancy}
\fancyhf{}
\fancyhead[L]{\Niveau}
\fancyhead[C]{\ChapitreCourt}
\fancyhead[R]{\Document}
\fancyfoot[L]{A. Fernandez}
\fancyfoot[R]{\thepage\ / \pageref{LastPage}}

% =====================================================
\begin{exercice}{1}
Une fonction affine s'écrit $f(x)=ax+b$.
\begin{tasks}(2)
\task $f(x)=4x+1$ : affine, $a=4$ et $b=1$.
\task $g(x)=-3x$ : affine, $a=-3$ et $b=0$. Elle est \textbf{linéaire}.
\task $h(x)=x^2-2$ : \textbf{pas affine}, à cause du $x^2$.
\task $k(x)=5-x=-x+5$ : affine, $a=-1$ et $b=5$.
\task $m(x)=6$ : affine, $a=0$ et $b=6$. Elle est \textbf{constante}.
\task $n(x)=2(x+3)=2x+6$ : affine, $a=2$ et $b=6$.
\end{tasks}
\end{exercice}

% =====================================================
\begin{exercice}{2}
\begin{enumerate}[label=\arabic*.]
\item $f(0)=-2\times 0+5=5$ \qquad $f(3)=-2\times 3+5=-6+5=-1$

$f(-1)=-2\times(-1)+5=2+5=7$ \qquad $f(1{,}5)=-2\times 1{,}5+5=-3+5=2$
\item $f(x)=1$\\
$-2x+5=1$\\
$-2x=1-5$\\
$-2x=-4$\\
$x=\dfrac{-4}{-2}=2$\\
L'antécédent de $1$ par $f$ est $2$.
\item $f(4)=-2\times 4+5=-8+5=-3$ : c'est l'ordonnée de $A$, donc $A$ \textbf{appartient} à la droite.

$f(-2)=-2\times(-2)+5=4+5=9\neq 8$, donc $B$ \textbf{n'appartient pas} à la droite.
\end{enumerate}
\end{exercice}

% =====================================================
\begin{exercice}{3}
\begin{minipage}[c]{0.5\linewidth}
Pour chaque droite, on calcule deux points.
\begin{enumerate}[label=\arabic*.]
\item \textcolor{blue}{$f_1(x)=2x-3$} : $f_1(0)=-3$ et $f_1(3)=3$.\\
Points $(0~;~-3)$ et $(3~;~3)$.
\item \textcolor{red}{$f_2(x)=-0{,}5x+2$} : $f_2(0)=2$ et $f_2(4)=0$.\\
Points $(0~;~2)$ et $(4~;~0)$.
\item \textcolor{green!50!black}{$f_3(x)=3$} : fonction constante, droite horizontale passant par $(0~;~3)$.
\end{enumerate}
\end{minipage}
\hfill
\begin{minipage}[c]{0.46\linewidth}
\centering
\begin{tikzpicture}[scale=0.5]
  \repere{-5}{6}{-4}{5}
  \draw[thick, blue, domain=-0.5:4] plot (\x,{2*\x-3}) node[right] {$f_1$};
  \draw[thick, red, domain=-5:6] plot (\x,{-0.5*\x+2}) node[below] {$f_2$};
  \draw[thick, green!50!black] (-5,3) -- (6,3) node[above left] {$f_3$};
  \foreach \p in {(0,-3),(3,3),(0,2),(4,0)} \fill \p circle[radius=3pt];
\end{tikzpicture}
\end{minipage}
\end{exercice}

% =====================================================
\begin{exercice}{4}
Méthode : $b$ est l'ordonnée du point d'intersection avec l'axe des ordonnées. Pour lire $a$, on avance de $1$ vers la droite et on compte de combien on monte (ou descend).

\begin{minipage}[c]{0.5\linewidth}
\begin{itemize}[label=$\bullet$]
\item $d_1$ : $b=1$ ; on monte de $2$, donc $a=2$.\\
$d_1$ : $y=2x+1$.
\item $d_2$ : $b=3$ ; on descend de $1$, donc $a=-1$.\\
$d_2$ : $y=-x+3$.
\item $d_3$ : $b=-2$ ; quand on avance de $2$, on monte de $1$, donc $a=\dfrac{1}{2}=0{,}5$.\\
$d_3$ : $y=0{,}5x-2$.
\item $d_4$ : droite horizontale, $a=0$.\\
$d_4$ : $y=-4$.
\end{itemize}
\end{minipage}
\hfill
\begin{minipage}[c]{0.46\linewidth}
\centering
\begin{tikzpicture}[scale=0.5]
  \repere{-5}{5}{-5}{6}
  \draw[thick, domain=-3:2.5] plot (\x,{2*\x+1}) node[right] {$d_1$};
  \draw[thick, domain=-3:5] plot (\x,{-\x+3});
  \node[above right] at (-3,6) {$d_2$};
  \draw[thick, domain=-5:5] plot (\x,{0.5*\x-2}) node[above] {$d_3$};
  \draw[thick] (-5,-4) -- (5,-4) node[above left] {$d_4$};
  % lectures de a
  \draw[blue, thick, ->] (-2,-3) -- (-1,-3);
  \draw[blue, thick, ->] (-1,-3) -- (-1,-1) node[midway, right] {\scriptsize $+2$};
  \draw[red, thick, ->] (0,3) -- (1,3);
  \draw[red, thick, ->] (1,3) -- (1,2) node[midway, right] {\scriptsize $-1$};
  \draw[green!50!black, thick, ->] (0,-2) -- (2,-2);
  \draw[green!50!black, thick, ->] (2,-2) -- (2,-1) node[midway, right] {\scriptsize $+1$};
\end{tikzpicture}
\end{minipage}
\end{exercice}

% =====================================================
\begin{exercice}{5}
On utilise $a=\dfrac{y_B-y_A}{x_B-x_A}$.
\begin{tasks}(2)
\task $a=\dfrac{8-2}{3-1}=\dfrac{6}{2}=3$
\task $a=\dfrac{1-5}{2-0}=\dfrac{-4}{2}=-2$
\end{tasks}
\end{exercice}

% =====================================================
\begin{exercice}{6}
On calcule $a=\dfrac{y_B-y_A}{x_B-x_A}$, puis on trouve $b$ avec les coordonnées de $A$ : $y_A=a\times x_A+b$.
\begin{tasks}(3)
\task $a=\dfrac{9-5}{3-1}=\dfrac{4}{2}=2$\\[0.3em]
$5=2\times 1+b$\\
$5=2+b$\\
$b=3$\\
$(AB)$ : $y=2x+3$
\task $a=\dfrac{-1-7}{2-(-2)}=\dfrac{-8}{4}=-2$\\[0.3em]
$7=-2\times(-2)+b$\\
$7=4+b$\\
$b=3$\\
$(AB)$ : $y=-2x+3$
\task $a=\dfrac{0-(-4)}{4-0}=\dfrac{4}{4}=1$\\[0.3em]
$A(0~;~-4)$ est sur l'axe des ordonnées, donc $b=-4$.\\
$(AB)$ : $y=x-4$
\end{tasks}
\end{exercice}

% =====================================================
\begin{exercice}{7}
Le sens de variation d'une fonction affine dépend du signe de $a$.
\begin{tasks}(2)
\task $a=5>0$ : $f$ est \textbf{croissante} sur $\mathbb{R}$.
\task $a=-2<0$ : $g$ est \textbf{décroissante} sur $\mathbb{R}$.
\task $a=0$ : $h$ est \textbf{constante} sur $\mathbb{R}$.
\task $a=-\dfrac{1}{2}<0$ : $k$ est \textbf{décroissante} sur $\mathbb{R}$.
\end{tasks}
\end{exercice}

% =====================================================
\begin{exercice}{8}
Méthode : on résout $f(x)=0$, puis on utilise le signe de $a$ : si $a>0$, la fonction est négative puis positive ; si $a<0$, elle est positive puis négative.

\medskip
\begin{minipage}[t]{0.48\linewidth}
1. $x-4=0 \iff x=4$ ; $a=1>0$.\par\smallskip
\centering\tabsignes{f}{4}{-,z,+}
\end{minipage}
\hfill
\begin{minipage}[t]{0.48\linewidth}
2. $2x+6=0 \iff 2x=-6 \iff x=-3$ ; $a=2>0$.\par\smallskip
\centering\tabsignes{g}{-3}{-,z,+}
\end{minipage}

\bigskip
\begin{minipage}[t]{0.48\linewidth}
3. $-3x+9=0 \iff -3x=-9 \iff x=3$ ; $a=-3<0$.\par\smallskip
\centering\tabsignes{h}{3}{+,z,-}
\end{minipage}
\hfill
\begin{minipage}[t]{0.48\linewidth}
4. $4-2x=0 \iff -2x=-4 \iff x=2$ ; $a=-2<0$.\par\smallskip
\centering\tabsignes{k}{2}{+,z,-}
\end{minipage}
\end{exercice}

% =====================================================
\begin{exercice}{9}
\begin{enumerate}[label=\arabic*.]
\item $P(x)=ax+b$ avec $P(4)=12$ et $P(10)=24$.

$a=\dfrac{24-12}{10-4}=\dfrac{12}{6}=2$.

$P(4)=12$ donc $2\times 4+b=12$, soit $8+b=12$, donc $b=4$.

$P(x)=2x+4$.
\item $a=2$ : chaque kilomètre coûte $2$~€.\\
$b=4$ : la prise en charge coûte $4$~€ (prix payé avant de rouler).
\item $P(15)=2\times 15+4=30+4=34$. Une course de $15$~km coûte $34$~€.
\item On résout $P(x)=30$ :\\
$2x+4=30$\\
$2x=26$\\
$x=13$\\
Avec $30$~€, on peut parcourir $13$~km.
\end{enumerate}
\end{exercice}

% =====================================================
\begin{exercice}{10}
\begin{enumerate}[label=\arabic*.]
\item Au départ, la charge est de $100\,\%$. Chaque heure, elle baisse de $8\,\%$. Après $t$ heures, elle a baissé de $8\times t$, donc $B(t)=100-8t=-8t+100$.

$B$ est une fonction affine avec $a=-8$ et $b=100$.
\item $a=-8<0$ donc $B$ est \textbf{décroissante}. C'était prévisible : la batterie se décharge au cours du temps.
\item On résout $B(t)=20$ :\\
$-8t+100=20$\\
$-8t=-80$\\
$t=\dfrac{-80}{-8}=10$\\
La batterie atteint $20\,\%$ au bout de $10$ heures.
\item $-8t+100=0$\\
$-8t=-100$\\
$t=\dfrac{-100}{-8}=12{,}5$\\
La batterie est vide au bout de $12{,}5$ heures, c'est-à-dire $12$~h~$30$~min.
\end{enumerate}
\end{exercice}

% =====================================================
\begin{exercice}{11}
\begin{minipage}[c]{0.52\linewidth}
\begin{enumerate}[label=\arabic*.]
\item \textcolor{blue}{$f(x)=12x$} : fonction \textbf{linéaire}.

\textcolor{red}{$g(x)=7x+60$} : fonction \textbf{affine}.
\item Pour $f$ : $f(0)=0$ et $f(20)=240$.

Pour $g$ : $g(0)=60$ et $g(20)=200$.
\item Les droites se coupent au point d'abscisse $12$ : pour $12$ entrées, les deux formules coûtent le même prix ($144$~€).
\end{enumerate}
\end{minipage}
\hfill
\begin{minipage}[c]{0.44\linewidth}
\centering
\begin{tikzpicture}[x=0.25cm, y=0.017cm]
  \draw[very thin, gray!40, xstep=1, ystep=20] (0,0) grid (24,300);
  \draw[thin, gray!70, xstep=4, ystep=100] (0,0) grid (24,300);
  \draw[->] (0,0) -- (25.5,0);
  \node[below=14pt] at (12,0) {\small nombre d'entrées};
  \draw[->] (0,0) -- (0,315) node[above] {\small prix (€)};
  \foreach \x in {4,8,...,24} \node[below] at (\x,0) {\scriptsize $\x$};
  \foreach \y in {100,200,300} \node[left] at (0,\y) {\scriptsize $\y$};
  \node[below left] at (0,0) {\scriptsize $0$};
  \draw[thick, blue] (0,0) -- (24,288) node[above left] {$f$};
  \draw[thick, red] (0,60) -- (24,228) node[below left] {$g$};
  \draw[dashed] (12,0) -- (12,144) -- (0,144);
  \fill (12,144) circle[radius=2pt];
\end{tikzpicture}
\end{minipage}

\begin{enumerate}[label=\arabic*., start=4]
\item $f(x)=g(x)$\\
$12x=7x+60$\\
$12x-7x=60$\\
$5x=60$\\
$x=12$\\
On retrouve $12$ entrées. Le prix est alors $f(12)=12\times 12=144$~€.
\item Après $x=12$, la droite de $g$ est en dessous de celle de $f$ : la formule B est moins chère. Elle est donc plus avantageuse à partir de $13$ entrées.
\item \textbf{Formule A :} $\dfrac{200}{12}\approx 16{,}7$, donc Léa peut faire $16$ entrées.

\textbf{Formule B :} on résout $7x+60\leqslant 200$, soit $7x\leqslant 140$, donc $x\leqslant 20$. Léa peut faire $20$ entrées.

Avec $200$~€, Léa fait plus d'entrées avec la \textbf{formule B}.
\end{enumerate}
\end{exercice}

\end{document}
